\documentclass[10pt,letterpaper]{article}
\usepackage{technical-testing}
\setDocTitle{Engine Agnostic Character Integration}
\setDocSubtitle{Implementation laboratory with reproducibility contract}
\setDocVersion{1.0}
\setDocStatus{Educational / Draft}
\setDocOwner{Game Development}
\setDocDate{2026-10-30}
\setDocSummary{Implementation laboratory with reproducibility contract}
\begin{document}
\makedoctitle
\makedocsummary

\section{Prerequisites and source crosswalk}
A 2D or 3D game engine; rig or sprite clips; frame telemetry. Book chapters: Chapters 8, 15--17, 73--77. Engineering procedure and code are supplementary.
\section{Input and baseline}
runtime architecture technical reference and state machine. Record shot ID, coordinate units, target fps, time convention, compiler/interpreter and fixture revision.
\section{Experiment procedure}
\begin{enumerate}
\item Import idle, run, attack and landing clips; define transition guards, contact markers and displacement owner.
\item Run a baseline and two controlled variations; change only one parameter per variation.
\item Capture source data (CSV or trace), a sampled frame strip, and test command output.
\item Identify at least one scenario where the intended visual effect conflicts with physical realism or performance limits.
\end{enumerate}
\section{Acceptance criteria}
Record animation trace, clip/rig compatibility, contact drift, 95th-percentile evaluation time and player interruption test results. No result is considered verified without stored output and test command.
\section{Required evidence bundle}
\begin{itemize}
\item Written hypothesis and expected perceptual result.
\item Environment/configuration manifest and exact execution commands.
\item Raw outputs, plots or contact sheets, comparison table and observed failure modes.
\item Engineering decisions, reviewers, and one prioritized revision.
\end{itemize}
\end{document}
